\begin{textsample}{POS Dim 1 – pb – Score 43.00 – pb/pb\_em\_38.txt}  \label{ex:f1_pos_126}
2.1.6.3.4 Processos \textbf{metodológicos} da Música Como a própria nomeação deste \textbf{tópico}, “ processos \textbf{metodológicos} ”, \textbf{aqui} propõe, a maneira de ensinar música é \textbf{um} processo \textbf{que} tem seus \textbf{métodos}, os quais devem estar em constante construção e com flexibilidade para as necessárias adequações, já \textbf{que} cada pessoa e cada turma é diferente e apresenta suas particularidades.
Desse modo, não é \textbf{interessante} \textbf{que} o/a professor/a faça sempre “ do seu jeito ” porque isso traduz inflexibilidade.
O mundo \textbf{atual} requer movimento, dinamismo, mudanças.
É \textbf{interessante}, sim, \textbf{que} possa seguir \textbf{um} ou mais \textbf{métodos}, ainda \textbf{que} os adapte a turmas e situações diferentes, pois, \textbf{certamente}, isso será necessário.
Dessa forma, vale \textbf{lembrar} \textbf{que} cada indivíduo tem sua maneira própria de aprender, portanto, o conjunto de indivíduos - a turma - vai somar também \textbf{uma} maneira própria para aprender.
Nesse sentido, o professor precisa encontrar a clave ( chave ) para depois descobrir o tom para começar a colocar na pauta o desenho do \textbf{método}, ou mesmo dos \textbf{métodos} \textbf{que} serão aplicados com aquela turma.
Além disso, na maioria das vezes, algumas metodologias são aplicadas numa mesma turma do ensino médio, sequencialmente, ou partes de algumas ao mesmo tempo, mesclando -as, mas, para \textbf{que} isso ocorra, é necessário \textbf{bastante} estudo e aperfeiçoamento formativo do/a professor/a.
Assim, não dá jamais para o professor querer ensinar música para seus jovens alunos/as exatamente como ele aprendeu.
Isso é \textbf{prova} de seu atraso pedagógico e rigidez na sua forma de \textbf{ver} o mundo.
Obviamente, isso acontece para aqueles \textbf{que} não têm a prática do estudo constante, \textbf{que} não buscam estar de acordo com o \textbf{que} a atualidade \textbf{exige} como, por exemplo, a aprendizagem com o uso de ferramentas digitais de maneira consciente, construtiva, criativa.
Dessa forma, destaca -se \textbf{que} os jovens estão em outro tempo, em outra “ vibe ”, como eles mesmos afirmam.
Já nasceram em meio a \textbf{um} mundo numa velocidade bem diferente daquele mundo \textbf{que} você, professor, encontrou quando tinha a mesma idade \textbf{que} eles/as se encontram agora. \textbf{Lembrando}, além disso, \textbf{que} utilizar novas tecnologias, aquelas sobre as quais os/as jovens estão imersos e, muitas vezes, até quase afogados, não é abandonar \textbf{métodos} mais usuais do ensino tradicional da música, como a voz, o corpo, a sucata, os sons da natureza, ou de instrumentos convencionais.
Pelo contrário, é encontrar a dosagem certa para cada turma.
Sendo assim, é essencial manter tradições como música em roda, cirandas, jogos e brincadeiras musicais, construção de instrumentos sonoros e muito do \textbf{que} já é encontrado no processo de ensino e aprendizagem musical.
Mas, também, deve ser utilizada a informática musical, como o recurso do computador e dos smartphones, por exemplo, para a criação, manipulação, execução, reprodução e veiculação da música, trabalhando a inclusão digital pelo acesso às novas tecnologias.
Mas, atenção, professor!
De nada adianta “ super equipar ” \textbf{uma} sala de aula com vários instrumentos e materiais sonoros de última geração sem mudar a forma de ensinar.
Não adianta adotar novas tecnologias e manter as mesmas metodologias. \textbf{Urge} acrescentar também \textbf{que} o ensino sobre o ritmo não pode se limitar ao \textbf{mero} pronunciar de \textbf{um} “ tá, tá, tá ” e repetições de células rítmicas.
Diferentemente, ele é a aproximação dos/as estudantes e seu sentimento de pertencimento ao \textbf{conseguirem} sentir sua respiração, seu caminhar, seu movimento e praticar a relação rítmica existente nos contos e cantos de versos e parlendas, seja de sua cultura local, seja de \textbf{um} repertório erudito, ou ainda de sua própria autoria.
Utilizar melodias, passagens rítmicas, textos, poemas e letras de canções sugeridas pelos/as próprios/as estudantes tende a motivá-los/as ao trabalho e à criatividade, principalmente com atividades lúdicas.
Nesse sentido, é possível adentrar no campo da notação musical.
Utilizar códigos a partir de \textbf{uma} criação simbólica, com signos gráficos para os sons, ruídos e mesmo para os tempos de silêncio, permite a compreensão de \textbf{que} a música pode ser representada graficamente.
Ao criar \textbf{uma} sequência sonora, esta pode ser recordada na aula seguinte, \textbf{uma} vez \textbf{que} foi grafada sem o uso de gravadores de som. \textbf{Veja} \textbf{aqui} o quanto é valioso deixar clara a importância de registrar graficamente a música, nesse caso, para a lembrança e o seu reuso.
Mais \textbf{uma} vez, \textbf{lembrando} \textbf{que} o/a estudante precisa compreender o porquê de estar aprendendo cada objeto de conhecimento.
A partir daí, \textbf{dependendo} da captação cognitiva da turma, passa -se a iniciar a leitura e a escrita de pequenos trechos, \textbf{simples}, em doses homeopáticas, sobre notação musical.
Se as aulas de Teoria Musical nas escolas regulares sofrem de muita resistência por parte dos estudantes, é, em grande parte, de responsabilidade do mau uso de metodologias.
Esse conteúdo chega com \textbf{uma} carga muito pesada, de maneira abrupta e sem conexão com a experiência individual deles. \textbf{Uma} metodologia muito bem aplicada no ensino de música no Brasil teve início no Rio de Janeiro com o professor Lucas Ciavatta, fundador do Instituto d’O Passo[21 ].
O \textbf{método} “ O Passo ” é \textbf{um} projeto especialmente democrático pela sua acessibilidade e propõe o ensino de Música por \textbf{intermédio} de \textbf{uma} vivência musical fundamentada nos princípios da inclusão e da autonomia, sustentadas pelos seus 4 pilares : corpo, imaginação, grupo e cultura.
É \textbf{um} \textbf{método} de alta eficácia e permite tranquilamente o trânsito ou a inserção de demais \textbf{métodos} na sua prática.
É, por assim \textbf{dizer}, \textbf{uma} maneira capaz de superar as \textbf{fronteiras} físicas de \textbf{um} instrumento ou da sala de aula convencionais.
No \textbf{método} “ Kodaly ”, do húngaro Zóltan Kodaly ( 1882-1967 ), propõe -se \textbf{que}, em jogos, brincadeiras, ou qualquer atividade lúdica, é possível trabalhar a escrita e a leitura musical, utilizando, por exemplo, as práticas de canto.
Heitor Villa-Lobos \textbf{inspirou} -se no \textbf{método} Kodaly, na criação do seu Canto Orfeônico, implantado nas escolas.
Outra referência no ensino de Música foi o austríaco Émile Jaques Dalcroze ( 1865-1950 ), \textbf{que} apresentou \textbf{um} \textbf{método} \textbf{que} \textbf{dizia} ser de desenvolvimento integral do ser humano, tanto física quanto intelectual, utilizando solfejos rítmicos, melódicos e improvisações.
A criação musical e a improvisação também são \textbf{basilares} no \textbf{método} “ Orff ”, do alemão Carl Orff ( 1895-1982 ), \textbf{que} desenvolveu em combinações de música e dança, atividades vocais e instrumentais em grupo.
Esses foram apenas alguns exemplos de dezenas de renomados \textbf{métodos} \textbf{que} devem ser bem estudados e explorados pelo professor de Música.
Conquanto \textbf{que} fique bem atento ao fato de \textbf{que} a sala de aula nas escolas não é \textbf{uma} academia para formação de virtuoses para exibições ao público.
É \textbf{um} ambiente de vivência musical em grupo, \textbf{que} pensa cada estudante como \textbf{um} \textbf{sujeito} de direito, \textbf{um} protagonista de sua formação cidadã.
É muito importante \textbf{que} toda a comunidade escolar tenha sempre a ciência de \textbf{que} \textbf{uma} apresentação artística é o resultado de \textbf{um} processo de vivências sociomusicais e jamais deve ser entendida como \textbf{um} “ show ” alheio ao transcurso desse processo.
Desta feita, as apresentações também devem ser estimuladas, mas não como a meta da aprendizagem.
A meta está no processo, e não no produto, muito embora também se aprenda \textbf{bastante} no momento de cada apresentação, pois elas ajudam na criatividade, desenvolvem emoções como a coragem, a força de \textbf{vontade} de buscar sempre o melhor de si mesmo, trabalha a coletividade, incentiva o movimento cultural na escola, envolvendo toda a comunidade, além, é claro, de valorizar o trabalho do aluno. \textbf{Uma} \textbf{tarefa} do professor \textbf{enquanto} construtor do seu planejamento contínuo é partir de \textbf{uma} \textbf{abordagem}, \textbf{uma} avaliação inicial diagnóstica, \textbf{que} pode ser, por exemplo, com a técnica de Roda de Conversa, entendendo o estágio natural de desenvolvimento daquela faixa etária, buscando \textbf{uma} captação de informações da vida musical \textbf{que} trazem dentro de si, como observa Keith Swanwick ( 2010, online ) : > O essencial é respeitar o estágio em \textbf{que} cada aluno se encontra.
Tendo isso em \textbf{mente}, é \textbf{preciso} seguir três princípios.
Primeiro, \textbf{preocupar} -se com a capacidade da criança de entender o \textbf{que} é proposto.
Depois, observar o \textbf{que} ela traz de sua realidade, as coisas com \textbf{que} também pode contribuir.
Por fim, tornar o ensino fluente, como se fosse \textbf{uma} conversa entre estudantes e professor.
Isso se faz muito mais demonstrando os sons do \textbf{que} com o uso de notações musicais[22 ] Swanwick ( 2003 ) prioriza e destaca a importância da experimentação sonora, seja em instrumentos ou em quaisquer materiais encontrados pelo estudante para \textbf{que} ele mesmo se sinta estimulado e vivenciado/a com a música do seu dia a dia, da sua cultura.
Logicamente, introduzindo novos conceitos à medida da receptividade da turma, bem como elementos, formas, estilos, padrões, repertório, entre outros.
Mas é \textbf{preciso} respeitar o universo sociocultural e afetivo do estudante.
Sabemos, \textbf{ademais}, \textbf{que} há muito de Jean Piaget e suas \textbf{teorias} de diversas fases da construção do conhecimento presente na obra de Keith Swanwick, tanto na “ Teoria Espiral do Desenvolvimento Musical ” como no “ Modelo Compreensivo da Experiência Musical ”, denominado C(L)A(S)P.
No entanto, não vamos nos ater \textbf{aqui} a Piaget, mas ao Modelo proposto por Swanwick.
Para ele, a aquisição do conhecimento musical vem a partir da vivência prática, complementada com conhecimentos técnicos e \textbf{literários} da música.
Nesse sentido, propomos a cada professor o estudo do \textbf{método} C(L)A(S)P, considerado \textbf{aqui} \textbf{uma} referência \textbf{que} pode atender as necessidades das carências de \textbf{um} modelo de ensino de Música do estado da Paraíba.
A sigla C(L)A(S)P representa os cinco parâmetros ou ações \textbf{que} o ser humano tem com a arte ( não são os parâmetros musicais como altura, duração etc ), na qual cada letra é a inicial desses parâmetros : - C de Composition ( composição ) ; - L de Literature Studies ( literatura sobre música ) ; - A de Appreciation ( apreciação ) ; - S de Skill Aquisition ( conhecimentos e habilidades sobre música ) e - P de Performance ( execução ).
A tradução mais adequada para nós é TECLA, em \textbf{que} : - T – Técnica ( manipulação de instrumentos, notação simbólica, audição ) ; - E – Execução ( cantar, tocar ) ; - C – Composição ( criação e improvisação ) ; - L – Literatura ( história da música e história de músicas ) e - A – Apreciação ( reconhecimento de estilos / forma / tonalidade / graus ).
Esses cinco parâmetros são equidistantes e equilibrados, ou seja, nenhum deles deve ser priorizado em detrimento de outros, haja vista \textbf{que}, se o professor percebe \textbf{que} a Apreciação está sendo mais utilizada do \textbf{que} a Execução, ele deve imediatamente contrabalançar suas ações.
Para ter êxito nesse processo, o professor deve cuidar para não cair nos antigos vícios de dar mais atenção a \textbf{uns} ( geralmente o \textbf{que} mais gosta ou domina ) do \textbf{que} aos outros.
Devem ser articulados, planejados e organizados de tal modo \textbf{que} as aulas não sejam apenas de audições de estilos musicais, ou \textbf{um} período inteiro de história da música, ou somente aulas de treinamento técnico de exercícios repetitivos e maçantes.
No \textbf{método} TECLA, de Swanwick, o estudante deve estar sempre se relacionando com a música, e não só com o conhecimento da música.
O professor deve, dessa forma, ter como ponto de apoio a tríade básica Apreciar, Executar e Compor.
Por outro \textbf{lado}, o/a estudante deve ser ensinado a pensar, refletindo nas suas decisões e compartilhando experiências com todo o grupo, aprendendo também a aceitar sugestões e novas ideias.
No parâmetro Apreciar ( A ), ele pode desenvolver \textbf{uma} audição crítica ; ao Executar ( E ), ele canta ou toca \textbf{um} instrumento convencional ou não, \textbf{que} pode até ser construído pelos próprios estudantes ou com o manuseio de computadores e smartphones ; ao Compor ( C ), ele pesquisa, explora, experimenta e cria suas músicas, podendo também gravar, editar, masterizar e compartilhar sua produção.
Entretanto, muitos ainda caem na cilada de trabalhar Literatura ( L ) como \textbf{uma} história da música isolada, sem \textbf{contextualização} com a Execução ( E ) e com a Apreciação ( A ).
Destacando \textbf{que} Literatura ( L ) e Técnica ( T ) são parâmetros tão importantes quanto a Tríade acima, no entanto, hierarquicamente, propõem -se construir bases \textbf{sólidas} para desenvolver as habilidades técnicas e os aprofundamentos teórico-históricos da música.
A PCEM/PB almeja para a Paraíba \textbf{um} ensino de Música como nunca foi praticado antes, com aulas \textbf{interessantes} e \textbf{que} motivem os alunos às benesses \textbf{que} traz essa linguagem.
O objetivo \textbf{central} da Educação Musical na Paraíba nas escolas não é formar musicistas ( se isso acontecer, então \textbf{que} ótimo também! ).
Mas seu objetivo é \textbf{que} com o professor de Música e as vivências musicais da escola, regidas por esse professor, os estudantes possam chegar à compreensão do \textbf{que} se passa nas músicas \textbf{que} escutam ou executam, pensar sobre o \textbf{que} estão \textbf{vivendo} e buscar sempre se \textbf{aprimorar}, individual e coletivamente. \textbf{Que} os estudantes de Música nas escolas da Paraíba possam se expressar e se satisfazerem musicalmente. [ 21 ] : Não temos \textbf{fronteiras}.
A ideia \textbf{central} é circular, estar em movimento. [... ].
O \textbf{que} fazemos – sejam cursos, aulas, oficinas, grupos musicais, shows, concertos ou eventos – está e estará sempre lá fora : nas ruas e nas praças, nas universidades e conservatórios, nas escolas, em museus ou presídios, ONGs ou empresas.
O Passo não se \textbf{restringe} a nenhum deles.
É para todos, onde for.
Assim como a música deve ser ( Extraído de https://www.institutodopasso.org/sobre -nos ). [ 22 ] : SWANWICK, Keith.
Keith Swanwick fala sobre o ensino de música nas escolas.
Entrevista concedida a Ana Gonzaga.
Nova Escola, online, em 01 de janeiro de 2010.

% matched lemmas: abordagem, ademais, aprimorar, aqui, atual, basilar, bastante, central, certamente, conseguir, contextualização, depender, dizer, enquanto, exigir, fronteira, inspirar, interessante, intermédio, lado, lembrar, literário, mente, mero, metodológico, método, preciso, preocupar, prova, que, restringir, simples, sujeito, sólido, tarefa, teoria, tópico, um, urgir, ver, viver, vontade
\end{textsample}
